\section{Verification purpose}

This case verifies convergence of Eulerian firebrand transport in a uniform one-dimensional wind. Surface propagation is suppressed and deposited firebrands ignite immediately, so the leading edge isolates generation, transport, deposition, and level-set ignition updates.

The test exposes grid- or Courant--Friedrichs--Lewy (CFL)-dependent advection speed, lost transient output, and inconsistent front extraction. Success requires the measured leading-edge rate to approach the prescribed wind speed over the accepted grid and CFL range.

\section{Mathematical formulation and numerical implementation}

\begin{sloppypar}
This sweep reconstructs the wind-Courant-number convergence experiment~\cite{Qin2025}.
Area-proportional firebrand production, Eulerian transport and accumulation,
and immediate ignition couple the arriving firebrands to the fire arrival time
and level-set field. Because surface spread is suppressed, sustained propagation
requires all these processes. The sweep changes only wind displacement per step
relative to cell size, with fixed physical domain and areal production. Errors
in transport, timestep scaling, immediate ignition, or arrival-time updates can
therefore appear as grid- or Courant-number-dependent leading-edge speed. The
wind-speed reference is an idealized numerical expectation, not a claim that
real firebrands travel deterministically at the ambient wind speed.

\end{sloppypar}

\section{Derivation of expected behavior}
With surface spread disabled and direct ignition enabled, the transported ember field is the sole propagation mechanism. Uniform advection therefore gives the analytical leading-edge trajectory $x_{LE}(t)=u_{wind}t$ and the normalized propagation rate $ROS/u_{wind}=1$. Grid and wind-CFL refinement should reduce numerical transport error without changing this limiting value.

% BEGIN EXPLICIT EXPECTED-RESULT DERIVATION
The generated timestep is not an independent free parameter:
\[
 \Delta t=CFL_w\frac{\Delta x}{6.71}.
\]
For the five required CFL values \(\{0.2,0.4,0.5,0.6,0.8\}\), the generated
timesteps are
\[
\begin{array}{c|ccccc}
\Delta x~(\mathrm{m})\backslash CFL_w&0.2&0.4&0.5&0.6&0.8\\\hline
5&0.149031&0.298063&0.372578&0.447094&0.596125\\
10&0.298063&0.596125&0.745156&0.894188&1.192250\\
20&0.596125&1.192250&1.490313&1.788376&2.384501\\
30&0.894188&1.788376&2.235469&2.682563&3.576751
\end{array}
\quad\mathrm{s}.
\]
These \(4\times5=20\) members form the acceptance set.  Uniform advection over
the 300--1500 m comparison interval predicts
\[
 \Delta t_{\mathrm{ref}}=\frac{1500-300}{6.71}=178.8376~\mathrm{s},
 \qquad \overline{ROS}_{\mathrm{ref}}
 =\frac{1200}{178.8376}=6.71~\mathrm{m\,s^{-1}}.
\]
At the 260-s stop the characteristic displacement is
\(6.71(260)=1744.6~\mathrm{m}\), plus the generated half-cell initial
coordinate where absolute position is needed.  Thus every required simulation condition
has the exact normalized reference
\(\overline{ROS}/U=1\) and \(E_U=|1-1|=0\); the implemented allowance is
\(E_U\leq0.10\).
% END EXPLICIT EXPECTED-RESULT DERIVATION

\section{Verification metrics and acceptance criteria}
ELMFIRE writes the level-set field \(\phi\) after every solver time step.
For every dump, the postprocessor reads the physical centerline, locates all
sign changes, and linearly interpolates between adjacent cell centers to find
the subcell locations where \(\phi=0\).  The most downwind zero crossing is
the leading edge \(x_{LE}(t)\).  Consecutive dumps give the stepwise rate
\begin{equation}
  ROS_i = \frac{x_{LE}(t_{i+1})-x_{LE}(t_i)}{t_{i+1}-t_i}.
\end{equation}
The reported average is the time-weighted mean of these stepwise rates while
both front observations lie between \SI{300}{m} and \SI{1500}{m}.  Thus every
solver step in the comparison interval contributes, and the result equals
the total leading-edge displacement divided by elapsed time.  This interval
excludes startup behavior, the maximum downwind flight scale near the source,
and the downstream boundary.  Raster shape, cell size, buffer exclusion,
dump-to-time matching, front-observation count, rate-interval count, and a
diagnostic least-squares coefficient of determination are recorded for every
result.  A CSV trajectory preserves each dump time, selected \(\phi\) raster,
leading-edge location, and stepwise rate of spread (ROS).

Every simulation condition with $CFL_{wind}\leq0.8$ passes when
\begin{equation}
  \left|ROS/u_{wind}-1\right|\leq 0.10.
\end{equation}
The higher-CFL points are retained to expose the onset of temporal
transport error but do not determine overall pass/fail status.

\subsection{Justification of metric quantification}

The required range \(CFL_{wind}\le0.8\) stays below the unit-CFL transport
limit and leaves a 20\% margin for the coupled source, transport, and direct-
ignition updates.  Thus the acceptance set probes the intended stable regime;
larger CFL values diagnose degradation without being used to redefine that
regime after results are known.  The 10\% speed band is deliberately wider
than the largest single-cell position uncertainty divided by the
\SI{1200}{m} comparison length, but much narrower than the order-one error
produced by a stalled front or by advancing trackers for only a fraction of a
solver step.  Time-weighting the rates over \SI{300}{m}--\SI{1500}{m} further
prevents an isolated grid crossing from controlling the decision.  Requiring
all 20 stable-regime simulation conditions guards against compensating resolution and
timestep errors; incomplete trajectories cannot establish convergence and are
therefore non-passing.

\subsection{Predeclared acceptance criteria}

Table~\ref{tab:acceptance-contract} summarizes the metrics fixed before
inspection of the calculated results. Numerical values and component statuses
are reported separately in Section~7.

\begin{table}[H]
\centering
\footnotesize
\caption{Predeclared verification metrics and acceptance criteria.}
\label{tab:acceptance-contract}
\begin{tabular}{@{}p{0.23\linewidth}p{0.47\linewidth}p{0.21\linewidth}@{}}
\toprule
Metric & Output, region, and aggregation & Acceptance criterion\\
\midrule
Required-output completeness & Transient level-set field rasters, dump times, and at least one valid front-rate interval between 300 and \SI{1500}{m}. & All 20 simulation conditions with \(CFL_{wind}\le0.8\) \\
Normalized mean-ROS error & \(E_U=|\overline{ROS}/u_{wind}-1|\), where \(\overline{ROS}\) is displacement divided by elapsed time over the comparison interval. & \(E_U\le0.10\) for every required simulation condition \\
Overall decision & All required simulation conditions must be complete and satisfy the normalized-speed limit; higher-CFL simulation conditions are diagnostic only. & PASS iff all required components pass \\
\bottomrule
\end{tabular}
\end{table}

\section{Simulation configuration and predicted outcome}

Figure~\ref{fig:input-configuration} records the prepared spatial inputs for a 10 m grid and wind Courant number 0.2.
These panels show the prepared spatial fields, remove the
two-cell numerical buffer, and retain map coordinates in metres.  The left panel shows
the most informative fuel or structure layout available for the case; the right panel
shows the cells selected by the non-positive initial level-set field, making a localized
ignition or firebrand-emission source explicit.

\begin{figure}[H]
\centering
\EvidenceFigure[width=0.98\textwidth,height=0.42\textheight,keepaspectratio]{../figures/input_configuration.pdf}
\caption{Whole-domain prepared input configuration for the representative simulation condition.  The
physical domain is shown without the numerical halo; coordinates, configured wind values,
fuel or structure layout, and initial ignition/source geometry are identified in the panels.}
\label{fig:input-configuration}
\end{figure}

The active physical domain is \SI{2400}{m} long by \SI{300}{m} wide.  Two
additional numerical cells are placed outside every physical boundary for
each resolution.  The initial burning cell is the first usable downwind cell
and the center row; it is therefore never in the buffer.  Fuel model 102,
flat terrain, a uniform \SI{15}{mph} wind from west to east, and deterministic
input rasters are used.  Crosswind dispersion and ember consumption are
disabled.

The current spotting interface uses the following explicit model selections:
\begin{itemize}
  \item generation: area-proportional;
  \item spotting distance: empirical;
  \item accumulation: Eulerian;
  \item ignition: immediate.
\end{itemize}
The customized downwind distribution has $\mu=2.18$, $\sigma=1.23$, and
$P_\varepsilon=0.01$.  Every ignited cell emits for \SI{10}{s}.  The configured areal generation rate is resolution-independent:
\begin{equation}
  GR_A=10\,[\mathrm{pcs\,m^{-2}\,s^{-1}}].
\end{equation}
ELMFIRE therefore emits \(N=GR_A\,\Delta x^2\,\Delta t\) firebrands from one burning square cell during one solver step. This is the appropriate source definition because surface-fire spread and its fireline-intensity-dependent power are disabled.

\subsection{Resolution sweep and predicted trend}
Four grid spacings, $\Delta x=5$, 10, 20, and \SI{30}{m}, are crossed with
wind-based CFL values 0.2, 0.4, 0.5, 0.6, 0.8, 1.0, and 1.2, for 28 simulation conditions.  Each
time step is prescribed by
\begin{equation}
  \Delta t=CFL_{wind}\frac{\Delta x}{u_{wind}}.
\end{equation}
Because timestep variation is the verification variable, these fractional values are not rounded. For each member preprocessing uses \(N=\lceil260/\Delta t\rceil\) and writes \(t_{stop}=N\Delta t\), recording the requested 260~s, aligned stop, and integer step count in the simulation specification. All simulation conditions retain the same physical domain; only the number of active and
buffer cells changes.

\section{Actual simulation results}

Figure~\ref{fig:domain-result} shows the representative transported level-set field from a 10 m grid and wind Courant number 0.2.
The displayed field is taken from the simulated results, masks missing values, and excludes
the two-cell numerical buffer.  The zero contour identifies the transported front, while the signed field exposes domain-wide numerical structure away from the extracted centerline.

\begin{figure}[H]
\centering
\EvidenceFigure[width=0.96\textwidth,height=0.50\textheight,keepaspectratio]{../figures/domain_result.pdf}
\caption{Whole-domain transported level-set field from the representative completed ELMFIRE run.  This
domain-scale result supplements, rather than replaces, the higher-level verification
profiles, convergence plots, ensemble statistics, and calculated acceptance metrics below.}
\label{fig:domain-result}
\end{figure}

Figure~\ref{fig:transport-convergence} first presents the complete wind-CFL
and grid-resolution sweep in the reference coordinate system. Each point is
the normalized speed calculated from a complete transient level-set history.
\begin{figure}[H]
  \centering
  \EvidenceFigure[width=0.95\linewidth,height=0.45\textheight,keepaspectratio]{../figures/eulerian_transport_convergence.pdf}
  \caption{Time-averaged Eulerian firebrand-driven spread speed normalized by
  the uniform wind speed as a function of wind-based CFL number. The shaded
  region marks the CFL range subject to the 10\% acceptance criterion.}
  \label{fig:transport-convergence}
\end{figure}

Figure~\ref{fig:front-extraction} is the supplemental diagnostic that traces
one representative point from the level-set contour through the reported mean.
It expands one member of the sweep,
\texttt{\Metric{representativevariant}}, to show the complete path from the
level-set field to the scalar convergence metric.  Panel (a) shows one actual
centerline profile, zoomed around the extracted front.  In this direct-ignition,
zero-surface-spread test, ELMFIRE's level set is nearly binary away from the
interface: burned and unburned cell centers remain close to $-1$ and $+1$.
This is expected and is not a corrupted field.  Adjacent cell-center values
$\phi_j$ and $\phi_{j+1}$ with opposite signs bracket the front.  Linear
interpolation gives
\begin{equation}
 x_{LE}=x_j-\phi_j\frac{x_{j+1}-x_j}{\phi_{j+1}-\phi_j}.
\end{equation}
All centerline sign changes are evaluated, and the most downwind crossing is
selected.  This matters after spotting because a line can contain several
separate burned regions and therefore several fire fronts.

Panel (b) assembles those independently extracted locations into
$x_{LE}(t)$.  The gray interval is the prescribed \SIrange{300}{1500}{m}
comparison window expressed in time.  The dashed least-squares line is a
trajectory-linearity diagnostic; its slope is
\Metric{representativefitslope}\,\si{m.s^{-1}}, but it is not the reported
acceptance metric.

Panels (c) and (d) show how that metric is calculated.  Panel (c) displays
a 20-second detail from the start of the averaging interval so the alternating
peak and trough values are individually visible.  Each pair of consecutive
dumps within the comparison window supplies
$ROS_i=\Delta x_{LE,i}/\Delta t_i$.  Because the final solver step can have a
different duration, the average is explicitly time weighted:
\begin{equation}
 \overline{ROS}
 =\frac{\sum_i ROS_i\Delta t_i}{\sum_i\Delta t_i}
 =\frac{x_{LE}(t_N)-x_{LE}(t_0)}{t_N-t_0}.
\end{equation}
The representative result is
$\overline{ROS}=\Metric{representativemeanros}\,\si{m.s^{-1}}$, or
$\overline{ROS}/u_{wind}=\Metric{representativenormalizedros}$.  Thus the
oscillatory instantaneous values caused by cell-by-cell front motion are not
mistaken for the effective transport speed.

\begin{figure}[H]
  \centering
  \EvidenceFigure[width=\linewidth,height=0.72\textheight,keepaspectratio]{../figures/representative\_front_extraction.pdf}
  \caption{Detailed extraction for the representative $\Delta x=\SI{10}{m}$,
  $CFL_{wind}=0.5$ simulation.  (a) Linear interpolation of centerline
  $\phi=0$ crossings; (b) resulting leading-edge trajectory and diagnostic
  fit; (c) a 20-second zoom of consecutive-dump finite-difference ROS and
  its time-weighted mean; (d) convergence of the cumulative normalized mean
  to its reported value.}
  \label{fig:front-extraction}
\end{figure}

\section{Calculated metrics and verification decision}
\begin{table}[htbp]
\centering
\begin{tabular}{lr}
\toprule
Quantity & Value \\
\midrule
Overall status & \Metric{status} \\
Prepared simulation conditions & \Metric{numvariants} \\
Simulation conditions with simulation results & \Metric{numvariantscomputed} \\
Required simulation conditions & \Metric{numrequiredvariants} \\
Required simulation conditions evaluated & \Metric{numrequiredvariantscomputed} \\
Reference wind speed [m/s] & \Metric{referencewindmps} \\
Maximum relative error & \Metric{maxrelativeerror} \\
Acceptance CFL limit & \Metric{acceptancecflmax} \\
\bottomrule
\end{tabular}
\caption{Verification status after processing all transient level-set outputs.}
\end{table}

\section{Technical interpretation}
The expanded sweep fails its aggregate criterion because the
$\Delta x=\SI{30}{m}$, $CFL_{wind}=0.5$ member gives
\[
  ROS/u_{wind}=\Metric{coarsecflofiveNormalizedros}
\]
corresponding to a
relative error of \Metric{coarsecflofiveRelativeerror}.  This is slightly
larger than the allowed 0.10.  The $\Delta x=\SI{10}{m}$ representative case
remains within tolerance and is used above to explain the extraction
procedure. All 28 input-matched simulation conditions completed and the generated controls match the declared wind, grid, and CFL sweep. The isolated coarse-grid excess is therefore a numerical transport/level-set discrepancy in the current source behavior, not an incomplete run or an admissible configuration correction.

\section{Scope and limitations}
Each prescribed condition is simulated and analyzed independently. The
comparison depends on complete results for the stated grid and timestep sweep,
not on whether a simulation or report preparation procedure finishes.

\begin{thebibliography}{9}
\bibitem{Qin2025}
Y. Qin, \emph{A Physics-Based Eulerian Framework for Modeling Firebrand
Showering in Regional-Scale Wildland and Wildland--Urban Interface Fire
Simulations}, Ph.D. dissertation, University of Maryland, College Park, 2025.
\end{thebibliography}
